% TOP_PROBLEM: {"id": 2302070, "title": "Research Problems in Function Theory — Problem 2.70", "category_id": 8, "category": {"id": 8, "name": "analysis", "display_name": "Analysis", "description": "Limits, continuity, calculus, and function theory.", "slug": "analysis", "order_index": 8, "created_at": "2026-07-31T15:26:25.671Z"}, "status": "open"}
% FIRST_POSTED: null
% ENTRY_KIND: statement_only
% Imported from: batch12.tex; UnsolvedMath ID 2302070
% Compile twice: lualatex 2302070.tex
\documentclass[11pt]{article}
\usepackage{fontspec}
\setmainfont{Latin Modern Roman}
\usepackage{amsmath,amssymb,mathtools,unicode-math}
\setmathfont{Latin Modern Math}
\usepackage{xurl,hyperref}
\hypersetup{colorlinks=true,urlcolor=blue,pdftitle={UnsolvedMath Problem 2302070}}
\newcommand{\sref}[2]{\hyperlink{source:#1:#2}{[S#2]}}
\newcommand{\cataloguescope}[1]{\paragraph{Mathematical question.}#1}
\begin{document}
% UnsolvedMath ID 2302070; source catalogue code AMR-022-2070
\setcounter{section}{2302069}
\section{Research Problems in Function Theory — Problem 2.70}\label{2302070}
{\small\sffamily UnsolvedMath ID 2302070 · Analysis}
\subsection{Definitions and mathematical statement}

\paragraph{Shared notation.}$\mathbb N=\{1,2,\ldots\}$, and $\mathbb Z,\mathbb Q,\mathbb R,\mathbb C$ denote the integers, rationals, reals and complex numbers. Any other conventions are specified locally.


For an entire function $f$, put $M(r,f)=\max_{|z|=r}|f(z)|$, $T(r,f)=(2\pi)^{-1}\int_0^{2\pi}\log^+|f(re^{i\theta})|\,d\theta$, and $\rho(f)=\limsup_{r\to\infty}\log\log M(r,f)/\log r$, where $\log^+t=\max(0,\log t)$. If $E$ is a nonzero entire function with nonzero zeros $a_j$, repeated according to multiplicity, define $\lambda(E)=\inf\{s>0:\sum_j|a_j|^{-s}<\infty\}$, with value $0$ for a finite zero set and $\infty$ when no such $s$ exists.
The source update reports counterexamples of Bergweiler and Eremenko for every order greater than $1/2$; this is the historical question, rather than an assertion that the proposed conclusion holds.
\paragraph{Statement recorded in the supplied source.}
Let $H$ be transcendental entire with $0<\rho(H)<\infty$ and $\rho(H)\notin\mathbb Z$, and let $f_1,f_2$ be linearly independent entire solutions of
\[w''(z)+H(z)w(z)=0.\]
Must $E=f_1f_2$ satisfy $\lambda(E)=\infty$? \sref{2302070}{2}
\subsection{Short English statement}
Must the product of two independent solutions have zeros of infinite exponent of convergence when the coefficient has finite nonintegral order?
\subsection{Sources}
\begin{description}
\item[\hypertarget{source:2302070:1}{S1}] ULAM.AI and the UnsolvedMath Contributors, UnsolvedMath: A Curated Collection of Open Mathematics Problems, record 2302070 (AMR-022-2070). CC BY 4.0. \url{https://huggingface.co/datasets/ulamai/UnsolvedMath}.
\item[\hypertarget{source:2302070:2}{S2}] W. K. Hayman and E. F. Lingham, \emph{Research Problems in Function Theory}, arXiv:1809.07200v2 (2018), Problem 2.70 and Update 2.70. \url{https://arxiv.org/abs/1809.07200}.
\end{description}
\label{end:2302070}
\end{document}
